\section{Scope, Limitations, and Falsifiability}

The compact-fiber derivation gives a structural electromagnetic bridge
coupling
\[
\alpha_F=\frac{1}{x_{\rm fiber}},
\]
where \(x_{\rm fiber}\) is the positive root of
\[
x^4
-
137x^3
-
\frac{74}{15}x^2
+
\frac{296}{105}
=
0.
\]
The claim is conditional: given the compact-fiber carrier, the
radiative/vibrational bridge-placement readout, the carrier-access
composition rule, the closure-budget accounting, the self-consistent
bridge-sharing rule, and the closed-return correction stated above, the
fixed-point equation and its positive root follow algebraically.

The result should therefore be read as a conditional structural
derivation of a low-energy electromagnetic bridge coupling, not as an
unconditional derivation from bare postulates alone and not as a completed
derivation of every empirical route by which \(\alpha\) is extracted.
Empirical determinations of the fine-structure constant involve
additional anchors and correction structures. For example,
\[
R_\infty=\frac{\alpha^2m_ec}{2h}
\]
contains the electromagnetic coupling, but also mass, action, and
unit-readout structure. Recoil, \(g-2\), spectroscopy, Josephson, and
quantum Hall determinations likewise infer \(\alpha\) through
route-specific theoretical and metrological pipelines.

The present paper separates these layers. The compact-fiber fixed-point
equation gives the structural bridge coupling \(\alpha_F\). The physical
role identification
\[
\alpha_F=\alpha_{\rm low-energy}
\]
rests on the \(\alpha_F^2\) role in leading atomic binding.
Route-specific metrological closure remains downstream; the general
compact-fiber readout and matching-correction layer is treated in
\cite{Wells2026FiniteReadout}.

\subsection{Coefficient status}

The fixed-point equation contains no continuously adjusted coefficient
and uses no measured value of \(\alpha\) as an algebraic input. Its
coefficients have named compact-fiber origins:
\[
137=128+9,
\qquad
\frac{74}{15}=5-\frac1{15},
\qquad
\frac{296}{105}=\left(\frac{74}{15}\right)\left(\frac47\right).
\]
Here
\[
128=|Z_4\times Z_4\times Z_8|,
\qquad
9=|V_{\rm rad}^{\rm bridge}|=3^2,
\]
\[
5=N_m+1,
\qquad
\frac1{15}=\frac1{N_m^2-1},
\qquad
\frac47=\frac{N_m}{N_e-1}.
\]
Thus the polynomial has a fully specified compact-fiber coefficient
structure under the stated bridge-readout and closure rules. The
empirical value of \(\alpha\) is not used to set the base count, the
feedback burden, or the closed-return coefficient.

This distinction is important. The paper does not claim that every
bridge-readout and closure rule used here is derived from bare postulates
within this manuscript. It claims that, once those rules are stated, the
alpha bridge equation follows without continuously adjusted coefficients
and without using the measured value of \(\alpha\) to set the algebraic
coefficients.

\subsection{Claim boundary}

The derivation depends on the following load-bearing ingredients:
\begin{enumerate}[label=(\arabic*)]
\item the compact-fiber carrier
\[
F=Z_4\times Z_4\times Z_8;
\]

\item the sign-status and placement-readout construction
\[
|V_{\rm rad}^{\rm bridge}|=3^2=9;
\]

\item carrier-level disjoint-union access counting for the
material/radiative bridge;

\item closure-budget accounting for the feedback burden
\[
K=\frac{74}{15};
\]

\item self-consistent transport sharing,
\[
x=B+\frac{K}{x};
\]

\item the source/complement return coefficient
\[
\frac{N_m}{N_e-1};
\]

\item the derivative self-response sign of the closed return.
\end{enumerate}

The result would require revision if any of these ingredients failed. In
particular, the derivation would require revision if \(F\) were not the
relevant material closure carrier, if
\(|V_{\rm rad}^{\rm bridge}|=3^2=9\) were not the correct
bridge-placement count, if the electromagnetic bridge required
product-state counting \(F\times V_{\rm rad}^{\rm bridge}\) rather than
disjoint-union access, or if the feedback burden \(K=74/15\) were not a
valid closure-budget term. It would also require revision if the
second-order return term
\[
-\frac{K N_m}{(N_e-1)x^3}
\]
were not the correct closed-return correction, or if another same-order
scalar bridge correction were admitted.

Finally, the low-energy role identification would fail if \(\alpha_F\)
did not play the \(\alpha^2\) role in leading atomic binding, or if
future measurement-chain analysis showed that operational extractions of
\(\alpha\) read a systematically different quantity than the
compact-fiber bridge coupling.

\subsection{Same-order and higher-order boundaries}

The second-order correction is conditionally unique within the current
scalar two-node bridge graph. A same-order competitor would require at
least one of the following:
\begin{enumerate}[label=(\arabic*)]
\item a second independent first-order feedback source;
\item a new carrier pool;
\item a preimage-resolved observable rather than scalar bridge readout;
\item a radiation-side feedback mechanism contributing to \(K\);
\item an independent third transport node;
\item a direct cubic mechanism not generated by the first-order feedback.
\end{enumerate}

No such mechanism is admitted in the present derivation. The uniqueness
claim is therefore relative to the stated scalar bridge graph and the
framework rules used in this paper. It is not a claim that all future
radiation-side feedback mechanisms are impossible.

The quartic equation includes the first closed return of the first-order
feedback burden. If repeated closed returns are admitted, each additional
round trip is suppressed by at least
\[
\frac{C}{x^2},
\qquad
C=\frac{N_m}{N_e-1}=\frac{4}{7}.
\]
For
\[
x\approx137.036,
\]
this suppression factor is approximately
\[
3.0\times10^{-5}.
\]
Thus higher-order repeated returns, if present, are operationally small
at the current comparison scale. A complete all-orders return theorem
remains open.

\subsection{Numerical comparison and route scope}

The compact-fiber point value lies within the CODATA 2022 experimental
uncertainty band for the inverse fine-structure constant. This comparison
is not a statistical validation of the framework, since the present paper
does not assign a theoretical uncertainty to the imported bridge-readout
rules, the closure-budget rules, the closed-return construction, or
possible omitted same-order mechanisms. It should therefore be read as a
consistency check, not as an error-budget closure.

Some empirical routes to \(\alpha\), especially Rydberg/recoil and
\(h/m\)-based routes, involve the action anchor
\[
h.
\]
The compact-fiber program treats the structural content of \(h\) as a
separate theorem path \cite{Wells2026PlanckGravity}. That path is
relevant to empirical route correspondence, but it is not part of the
alpha fixed-point theorem. The
alpha theorem supplies the compact-fiber bridge coupling \(\alpha_F\);
the Planck/action theorem supplies separate structural action content;
route analysis then explains how empirical determinations combine
electromagnetic coupling, action, mass, charge, and readout conventions.

The Planck/action result is used here only as context for legacy-route
correspondence.

\subsection{Falsifiability and summary of scope}

The present paper supplies an internal structural prediction and a
comparison with the empirical low-energy coupling. It does not claim a
complete independent falsification program. The internal derivation is
falsifiable at the framework level by failure of any stated load-bearing
rule: carrier choice, bridge-placement count, disjoint-union access,
closure-budget accounting, self-consistent sharing, return coefficient,
or derivative self-response sign.

A stronger external test would require applying the same bridge-readout
and closure rules to a second observable without introducing new
structural choices. Such a test would distinguish a reusable compact-fiber
rule set from an isolated alpha-specific construction. The present paper
identifies the alpha bridge prediction and its failure modes, while
leaving broader route-by-route and second-observable closure to
subsequent work.

Within the stated framework boundaries, the compact-fiber derivation gives
\[
\boxed{
\alpha_F^{-1}=x_{\rm fiber}.
}
\]
The result is not merely numerical: each coefficient in the fixed-point
equation has a named structural origin. It is also not a completed
metrological derivation: route-specific empirical extraction pipelines
remain downstream theorem programs. The strongest claim of this paper is
therefore a conditional structural prediction for the low-energy
electromagnetic bridge coupling, with explicit route-readout,
theoretical-uncertainty, and independent-test boundaries.